\chapter*{プロローグ}
\epigraph{The street finds its own uses for things.\\(ストリートは、物事に対して独自の使い方を見つけ出す。)}{William Gibson, Burning Chrome}
リチャード・バンドラー博士（Dr.~Richard Bandler）に初めて会ったのは、もう10年も前のことだ。
彼に会えば自分は変われるかもしれないと大きな期待に胸を膨らませて、日韓NLPディレクターのアンナ・スィルに連れられ、成田空港から旅立ったのを今でも鮮明に覚えている。

当時の私は、子供の頃に母からプレゼントされたTIMEXの時計を愛用していた。ストップウォッチやラップ計測など機能が豊富で、物理学修士である私にとってはお気に入りの相棒だった。
ところが、成田空港で買い物をしている最中に不意にバンドが壊れ、床に落として壊してしまったのだ。少し落胆した気分のまま、私はアメリカ行きの飛行機に乗り込ったのである。

トレーナートレーニングの会場は、フロリダのフォーポイント・シェラトン。そびえ立つタワーホテルの一室が私の部屋だった。一階には広大なセミナールームがあり、プールやジムも併設されていた。

そして、トレーニングの幕が開いた。
大音響の音楽が鳴り響き、私たちはリズミカルな拍手で迎えられた。バンドラーのウィットに富んだジョークに会場は大爆笑に包まれ、誰もがその場にグイグイと引き込まれてった。

そのとき突然、バンドラーは客席を見回してこう言い放った。
\begin{quote}
「私はTIMEXみたいな安物の時計はしない。こういう高級時計を使うんだ。だが高級時計っていうのは、すぐに狂っちまう。だからしょっちゅうメンテナンスをし続けなきゃならないんだ」
\end{quote}

この瞬間、私の思考は完全に停止し、一気に深い催眠状態へと引きずり込まれたのだ。

それから10年、催眠とは、そして人間の認識とは何なのだろうと問い続けてきた。そして今、私はひとつの明確な答えにたどり着いたと感じている。

バンドラー博士には、数々の奇妙でパワフルな信念をインストールしていただいた。そのおかげで何度も決死の大冒険へ放り出され、ひどい目に遭う羽目にもなったが、そのすべてに感謝している。そして、ここに提示する真理へと至る無数のヒントを与えてくださったことに、心より深く感謝申し上げる。

\section*{警告：本書を開く前のプロトコル}
本書は、いわゆる「魔術」の完全な種明かしの書である。

もしあなたが、エゴが描く波乱万丈の大冒険や、胸を引き裂くような愛のドラマをまだ十分に味わい尽くしていないのなら、今すぐ本書を閉じるべきだ。いますぐリチャード・バンドラーの元へ駆け込み、脳のレジスタに愉快でデタラメで、とびきりパワフルな信念をインストールしてもらってくるがいい。この世界という壮大なフィクションを全力で楽しむことこそが、いまのあなたに必要なプロセスだからだ。

本書を開いてよいのは、この現実というゲームを骨の髄まで遊び倒し、もはや舞台の裏側――奈落に張り巡らされた配線と、照明を制御する冷徹な物理コード――を直視する覚悟ができた者だけである。

ひとつだけ警告しておく。
中途半端な覚悟で本書のアルゴリズムをなぞれば、あまりの真実味にめまいを覚え、すべての意味が蒸発したかのような深いニヒリズムの泥沼に足を取られるかもしれない。

しかし、もしあなたがその荒野を越えて先へ辿り着くならば、約束しよう。
そこには、神話も錯覚も脱ぎ捨てた者だけが足を踏み入れることのできる、かつて誰も経験したことのない、新次元の冒険が待っている。


